\section{Introduction}

CIFAR-10~\citep{krizhevsky2009} is a small-image benchmark that modern architectures have largely
saturated: with tens of millions of parameters and long schedules, the remaining
headroom is small. The problem changes when the model must be trained from
scratch under a hard budget of five million trainable parameters, and when the
evaluation set is hidden and may not match the training distribution. Both
constraints were active in this competition. The first forces careful
architecture and optimization choices; the second makes validation accuracy an
imperfect guide and turns the leaderboard into a noisy, partially observed
signal.

We approached the problem as a two-track collaboration. One track pushed clean
accuracy through architecture search, long schedules, and knowledge distillation
(KD)~\citep{hinton2015distilling} from large teachers trained on the same data. The other track studied
robustness: it built corrupted and photometric proxy validation sets from the
training split and ran controlled augmentation ablations. The two tracks shared
one codebase, one fixed 5,000-image holdout, and one set of conventions; all numbers use
the same holdout, and the TTA/EMA protocol is stated wherever it differs.

The main result is a 4.88M-parameter student that reaches 97.84\% with
horizontal-flip test-time augmentation (TTA) and 97.92\% with 10-view TTA. The
team's selected submissions scored 0.911 and 0.904 on the public leaderboard;
the private score was 0.880, placing 13th of 81 teams. The roughly ten-point gap between our
best holdout number and the private score is the central unsolved problem of the
project. We do not have the final evaluation labels, and we do not analyze the
provided unlabeled images here, so we do not tune against the gap or claim to
explain it.

This paper makes four contributions.
\begin{itemize}
\item A strong, reproducible 4.88M-parameter ResNet~\citep{he2016deep} recipe: long cosine training~\citep{loshchilov2017sgdr}, SAM~\citep{foret2021sam},
EMA, TrivialAugment~\citep{muller2021trivial}, Cutout~\citep{devries2017cutout}, mixup/CutMix~\citep{zhang2018mixup,yun2019cutmix}, and three-teacher KD. It improves a
flip-TTA baseline from 95.42\% to 97.84\% (97.92\% with 10-view TTA) and yields
the team's 0.911 public / 0.880 private submissions.
\item A controlled account of which ingredients mattered. The training recipe gave the largest gain and KD a smaller, schedule-confounded one; architecture changes within the 5M budget gave small ones;
several plausible tricks did not pay off.
\item A clean-neutral augmentation result: one-view AugMix~\citep{hendrycks2020augmix} leaves clean accuracy
unchanged but improves a composite corruption proxy by 2.2pp and a
photographic-jitter proxy by 1.9pp, replicated across two non-KD runs, with preliminary checkpoint evaluations (epoch 200 of 300) suggesting similar gains under
distillation. A JSD variant did not preserve clean accuracy and was rejected.
\item A negative-results record. We document failed techniques with paired
statistics and validation curves: augmentation policies, late continuations,
fresh 300-epoch sweeps, test-time adaptation, TTA variants, model splitting, and
seed ensembling.
\end{itemize}
